%%%%%%%%%%%%%%%%%%%%%%%%%%%%%%%%%%%%%%%%%%%%%%%%%%%%%%%%%%%%%%%%%%%%%%%%
% TAILORING GUIDE FOR FUTURE CHATGPT CLIENTS
%
% Purpose:
% This master CV provides modular content with multiple specialization options.
% Follow the instructions below to generate a position-specific CV.
%
% Tailoring Rules:
% 1. Review the tailoring configurations below to select relevant content.
% 2. After selection, REMOVE the corresponding subsection titles such as:
%       \textbf{Material Simulation and Analysis}
%       \textbf{Software Development and Automation}
%    Retain only the higher-level section titles, such as:
%       \textbf{Postdoctoral Researcher}, University of South Florida \hfill \textit{2022 – Present} \\
%
% Suggested Tailoring Configurations:
%
% 1. Software R&D / Scientific Computing Positions
%    - Keep:
%      * Professional Summary
%      * Technical Skills
%      * Software Development and Automation (Postdoc)
%      * High-Performance Computing Optimization (Postdoc)
%      * Debugging, Profiling, and Performance Tuning (Postdoc)
%      * Software Development and GPU Acceleration (PhD)
%      * HPC and Workflow Optimization (PhD)
%    - Optional:
%      * Scientific Communication and Collaboration (PhD)
%    - Remove:
%      * Materials Simulation and Analysis (Postdoc)
%      * Structure–Property Analysis (PhD)
%
% 2. Materials Scientist / Simulation Researcher Positions
%    - Keep:
%      * Professional Summary
%      * Technical Skills
%      * Material Simulation and Analysis (Postdoc)
%      * Structure–Property Analysis (PhD)
%      * HPC Optimization and Workflow Acceleration (Postdoc)
%      * Scientific Communication and Collaboration (PhD)
%    - Optional:
%      * Software Development and Automation (Postdoc)
%    - Remove:
%      * Debugging, Profiling, and Performance Tuning (Postdoc)
%      * Software Development and GPU Acceleration (PhD)
%
% 3. Multidisciplinary Leadership / Team Management Roles
%    - Keep:
%      * Professional Summary
%      * Technical Skills
%      * Leadership, Mentorship, and Community Engagement (Postdoc)
%      * Scientific Communication and Collaboration (PhD)
%      * Software and HPC sections (if leadership is technical)
%    - Optional:
%      * All technical sections based on role expectations
%    - Remove:
%      * None, unless character limit requires further trimming
%%%%%%%%%%%%%%%%%%%%%%%%%%%%%%%%%%%%%%%%%%%%%%%%%%%%%%%%%%%%%%%%%%%%%%%%
\documentclass[letterpaper,12pt]{article}
\usepackage[margin=1.00in, top=0.7in, headsep=20pt]{geometry}
\usepackage{titlesec}
\usepackage{enumitem}
\usepackage{hyperref}
\usepackage{parskip}
\usepackage{fontawesome}

\hypersetup{
  colorlinks=true,
  urlcolor=blue,
}

% Font
\usepackage[T1]{fontenc}
\usepackage[sc]{mathpazo}

\usepackage[backend=biber,sorting=mysort,style=numeric-comp,defernumbers,maxbibnames=20]{biblatex}
\addbibresource[label=articles]{articles.bib}
\titleformat{\section}{\large\bfseries}{}{0em}{}[\titlerule]
 % fix header
 \defbibheading{fix}[\refname]{\section*{#1}}
 

\AtBeginBibliography{\small}

  % This is the sorting scheme for the CV
  \DeclareSortingScheme{mysort}{
    \sort[direction=descending]{
      \field{year}
    }
    \sort[direction=descending]{
      \field[padside=left,padwidth=2,padchar=0]{month}
      \literal{00}
    }
  }

  % Use this for reversing the numbers (count down) in the bibliography
  % Count total number of entries in each refsection
  \AtDataInput{%
  \csnumgdef{entrycount:\therefsection}{%
  \csuse{entrycount:\therefsection}+1}}

  % Print the labelnumber as the total number of entries in the
  % current refsection, minus the actual labelnumber, plus one
  \DeclareFieldFormat{labelnumber}{\mkbibdesc{#1}}    
  \newrobustcmd*{\mkbibdesc}[1]{%
  \number\numexpr\csuse{entrycount:\therefsection}+1-#1\relax}

% Custom sections
\usepackage[]{titlesec}
\titleformat{\section}{
  \rmfamily\mdseries\Large\bfseries
  }{}{}{}[\vspace{0.1ex}\titlerule]
\titlespacing*{\section}{0ex}{3ex}{0.8ex}
\titlespacing*{\therefsection}{0ex}{3ex}{0.8ex}
\titlespacing*{\refsection}{0ex}{3ex}{0.8ex}

\newlist{tabitemize}{itemize}{1}
%\setlist[tabitemize]{label=\textbullet,nosep,align=parleft,leftmargin=*,}
\setlist[tabitemize]{label=\textbullet,itemsep=1pt,align=parleft,leftmargin=*,}

%% Customize page headers
%\pagestyle{myheadings}
%\markright{\name}
%\thispagestyle{empty}
\usepackage[]{fancyhdr}
\pagestyle{fancy}
\fancyhead{}
\renewcommand{\headrulewidth}{0pt}
%\fancyhead[RO,LE]{Pierre Kawak, Ph.D.}
\fancyfoot{}
\fancyfoot[LO,CE]{\thepage}
\fancyfoot[RO,LE]{Pierre Kawak, Ph.D.}
\begin{document}

\begin{center}
  {\LARGE \textbf{Pierre Kawak, Ph.D.}}\\ \vspace{2mm}
  Open to relocation \quad Eligible to work in the U.S. \quad No sponsorship needed \\ \vspace{1.5mm}
  \faPhone\ (801) 762-7999 \quad \faEnvelope\ \href{mailto:pskawak@gmail.com}{pskawak@gmail.com} \quad \faLink\ \href{https://linktr.ee/pkawak}{linktr.ee/pkawak}
\end{center}

\vspace{-0.3\baselineskip}
\section*{Professional Summary}
Computational polymer physicist with expertise in molecular simulation, statistical mechanics, and polymer theory. Extensive experience leading collaborative projects at the intersection of materials science, chemical engineering, and polymer physics.
\vspace{0.5em}
\begin{tabitemize}[leftmargin=*]
    \item Expert in developing, optimizing, and deploying molecular dynamics and Monte Carlo simulations to uncover structure–property relationships in polymers and soft materials.
    \item Experienced in leading cross-functional teams, managing up to 11 researchers, and mentoring graduate and undergraduate students on computational methods and scientific communication.
    \item Proven ability to architect and accelerate analysis pipelines for polymer simulations, including profiling and optimizing in-house codes in Python, C++, and CUDA.
    \item Demonstrated success translating molecular insights into design principles for tough, recyclable, and sustainable materials, including elastomeric nanocomposites and copolymer systems.
    \item Strong record of scientific communication with multiple invited talks, publications in top journals, and community leadership in the polymer physics community.
    \item Passionate about building inclusive scientific communities, with leadership experience organizing professional development events, symposia, and mentoring programs.
\end{tabitemize}

\vspace{-0.3\baselineskip}
\section*{Technical Skills}
\begin{tabitemize}[leftmargin=*]
    \item \textbf{Simulation and Modeling:} Molecular dynamics (MD), LAMMPS, GROMACS, HOOMD-blue, Monte Carlo (MC), advanced sampling algorithms, coarse-grained modeling, atomistic modeling, OPLS, non-equilibrium simulations, phase behavior analysis, free energy calculations.
    \item \textbf{Software Development:} Architecting, profiling, and deploying scientific software in Python, C, C++, CUDA, Bash, Julia, R, Matlab, and FORTRAN; experience with collaborative version control, software development, unit testing, and continuous integration and development using Git.
    \item \textbf{Analysis and Data Processing:} Custom analysis pipelines for molecular simulation data using Python (NumPy, SciPy, Matplotlib, Pandas, scikit-learn), Bayesian optimization, statistical modeling, bash scripting, and parallel computation (CUDA, MPI, SLURM).
    \item \textbf{High-Performance Computing:} Experience with HPC job scheduling (SLURM), distributed computing (MPI), running simulations on national supercomputing facilities (XSEDE, ACCESS), and GPU-accelerated computing (CUDA, Nsight, nvprof, nvidia-smi).
    \item \textbf{Debugging and Profiling:} GDB, Nvprof, nvidia-smi, Nsight, perf, valgrind, callgrind.
    \item \textbf{File Management:} Automated backups, compression, terabyte-scale datasets.
    \item \textbf{Data Visualization:} Scientific visualization with Matplotlib, Inkscape, VMD, and OVITO.
    \item \textbf{Process Simulation and Thermodynamics:} Experience with process simulation tools (Aspen HYSYS) and instruction in thermodynamic modeling.
    \item \textbf{Leadership and Project Management:} Leading cross-functional teams, mentoring junior researchers, organizing professional development workshops and symposia.
    \item \textbf{Scientific Communication:} LaTeX, Scientific writing (6 publications), Technical presentations (27+ Talks), Mentorship
\end{tabitemize}

\vspace{-0.3\baselineskip}
\section*{Research Experience}

\textbf{Postdoctoral Researcher}, University of South Florida \hfill \textit{2022 – Present} \\
\textit{Advisor: Prof.~David Simmons}
\vspace{0.5em}

\textbf{Material Simulation and Analysis}
\begin{tabitemize}[leftmargin=*]
    \item Developed molecular dynamics simulations to investigate stress relaxation, deformation, and nanoscale reinforcement in polymer nanocomposites.
    \item Simulated thermo-elastoplastic deformation and rheology of filled rubber, identifying nanoscale toughening mechanisms.
    \item Automated selection, generation, and testing of atomistic sequence-designed macromolecules for enhanced thermal properties.
    \item Identified copolymer sequences with enhanced glass transition temperatures without altering processing or feedstock chemistry.
    \item Applied Bayesian optimization to fit relaxation dynamics and extract glass transition behavior.
    \item Led multi-dimensional parameter sweeps (nanoparticle size, chemistry, volume fraction) to optimize nanocomposite designs.
    \item Collaborated with experimentalists to validate molecular predictions against thermal and mechanical performance targets.
    \item Designed statistical models and applied machine learning to extract insights from large molecular datasets.
    \item Developed methods to extract stress hotspots from 3D trajectories for identifying mechanisms of reinforcement.
\end{tabitemize}

\vspace{0.5em}
\textbf{Software Development and Automation}
\begin{tabitemize}[leftmargin=*]
    \item Architected Python/C++ simulation and analysis tools to enable high-throughput, physics-informed property prediction.
    \item Developed modular, reusable rheology, thermal, and stress analysis toolkits integrated into team-wide simulation workflows.
    \item Automated data analysis pipelines handling 100+ TB datasets, reducing turnaround time by 90\%.
    \item Built bash, Slurm, and MPI automation for large-scale HPC job orchestration, checkpointing, and storage.
    \item Documented group codebases and created tutorials for reproducible scientific computing.
    \item Modernized a 50k+ line C++ molecular analysis toolkit (AMDAT): consolidated fragmented local versions into a single GitHub repo, eliminated segfaults and compiler warnings, and introduced unit tests and CI for reproducible releases.
\end{tabitemize}

\vspace{0.5em}
\textbf{High-Performance Computing Optimization and Workflow Acceleration}
\begin{tabitemize}[leftmargin=*]
    \item Led optimization of 50+ TB data pipelines across HPC clusters, improving data throughput by 90\%.
    \item Architected scalable storage and compression pipelines for terabyte-scale datasets; improved data access and reproducibility across research projects.
    \item Streamlined multi-node CPU and GPU workflows for molecular simulations.
    \item Reduced simulation runtime by 90\% through parallelization and high-performance computing resource optimization.
    \item Awarded NSF ACCESS Compute Resource Grant for advancing scalable HPC-enabled molecular simulations.
\end{tabitemize}

\vspace{0.5em}
\textbf{Debugging, Profiling, and Performance Tuning}
\begin{tabitemize}[leftmargin=*]
    \item Debugged and profiled C++/Python group codebases using GDB, valgrind, and profiling tools (nsight, nvprof), resolving memory leaks and runtime bottlenecks.
    \item Improved software modularity, runtime, and scalability for multi-month, high-dimensional simulations.
    \item Led efforts to standardize code maintainability and scientific reproducibility through internal documentation and technical presentations.
\end{tabitemize}

\vspace{0.5em}
\textbf{Leadership, Mentorship, and Community Engagement}
\begin{tabitemize}[leftmargin=*]
    \item Led and mentored 11+ researchers in HPC, molecular simulation, version control, and scientific computing practices.
    \item Designed structured onboarding programs, documentation, and tutorials to train new users.
    \item Organized professional development workshops, DEI programs, and research symposia for 200+ researchers across multiple organizations.
    \item Presented research at 17+ institutional, industrial, and academic conferences, earning multiple awards (GRC 2024, USF 2023).
    \item Delivered 17 conference talks; earned awards at GRC (2024) and USF Symp.~(2023).
    \item Recognized with the APS Career Mentor Fellowship for leadership in mentoring and community building.
  \item Refactored two legacy C++ toolkits into modular, open-source frameworks, enhancing community adoption and productivity across simulation \& analysis workflows.
  \item Transformed two legacy simulation and analysis codebases—initially underused and fragmented—into robust, modular Python/C++ frameworks, significantly improving lab-wide productivity and enabling their open-source release to support the broader polymer simulation community.
  \item Open-sourced AMDAT (Amorphous Molecular Dynamics Analysis Toolkit) on GitHub, adding user and developer docs, contribution guidelines, and examples that enabled adoption by external soft-matter simulation groups.
  \item Led authorship of an AMDAT software manuscript and coordinated invited talks to drive uptake, framing the toolkit as standard infrastructure for amorphous/polymer MD analysis.
\end{tabitemize}

\textbf{Ph.D.~Researcher}, Brigham Young University \hfill \textit{2017 – 2022} \\
\textit{Advisor: Prof.~Douglas Tree}

\vspace{0.5em}
\textbf{Software Development and GPU Acceleration}
\begin{tabitemize}[leftmargin=*]
  \item Built two GPU-accelerated Monte Carlo engines in C++/CUDA ($\sim$ 80,000 lines) to simulate polymer crystallization and free energy landscapes.
    \item Developed advanced sampling algorithms (Configurational Bias Monte Carlo, Wang-Landau) to efficiently explore high-dimensional molecular states.
    \item Solved critical convergence failures in polymer crystallization simulations by architecting a scalable, GPU-accelerated multi-window Monte Carlo framework — cutting runtimes from weeks to hours and enabling previously impossible sampling.
    \item Designed multi-window, multi-walker sampling algorithms to enhance convergence in high-dimensional MC simulations of polymer melts, enabling tractable modeling of metastable and glassy states.
    \item Designed GPU-optimized data structures (cell lists) and pairwise interaction search algorithms, achieving significant performance improvements.
    \item Utilized NVIDIA profiling tools (nsight, nvprof, nvidia-smi) to debug, optimize, and scale GPU-accelerated codes.
\end{tabitemize}

\vspace{0.5em}
\textbf{Structure–Property Analysis and Data Processing}
\begin{tabitemize}[leftmargin=*]
    \item Developed numerical algorithms to quantify nanoscale crystalline and orientational order from 3D simulation trajectories.
    \item Processed large-scale simulation data to generate thermodynamic observables, including free energy derivatives and heat capacity profiles.
    \item Extracted phase diagrams and structure–property relationships through systematic parameter sweeps and automated analysis pipelines.
\end{tabitemize}

\vspace{0.5em}
\textbf{High-Performance Computing and Workflow Optimization}
\begin{tabitemize}[leftmargin=*]
    \item Leveraged HPC resources to perform long-timescale, high-throughput Monte Carlo simulations using GPU acceleration.
    \item Automated simulation and data analysis workflows to enable reproducible exploration of high-dimensional parameter spaces.
    \item Developed a parallel sampling engine with 320+ concurrent jobs using multi-walker coordination across CPUs and GPUs — dramatically improving ergodicity and supporting multiple publications and an NSF CAREER award.
    \item Led development of a multi-walker, multi-window parallel sampling architecture deployed across 320+ concurrent jobs, integrating CPU/GPU scheduling and inter-process communication to boost sampling ergodicity and simulation throughput.
    \item Implemented fail-safe heartbeat system for resilient simulation orchestration and dynamic restart logic in long-timescale jobs.
    \item Built an automated, robust, and modular Python workflow system with JSON-based simulation state tracking, adaptive resource allocation, and dynamic error handling across SLURM-managed HPC jobs.
    \item Executed large multidimensional parameter sweeps for coarse-grained polymers.
\end{tabitemize}

\vspace{0.5em}
\textbf{Scientific Communication and Collaboration}
\begin{tabitemize}[leftmargin=*]
  \item Presented research findings at academic conferences, advancing understanding of nanoscale phenomena in polymer materials.
  \item Mentored 4 undergraduates; co-authored 2 papers and 6 conference abstracts.%; supported grad school admissions.
  \item Won APS Distinguished Student Award and BYU Research Presentation Award.
  \item Contributed key data to a successful \$500K NSF CAREER proposal.
  \item Simulation framework became foundational infrastructure for lab-wide research and supported a successful NSF CAREER proposal on polymer crystallization.
  \item Conceived and drove 3 first-author manuscripts from idea to submission in 2–6 months, coordinating simulation campaigns, replication, and iteration with collaborators.
\end{tabitemize}

\textbf{Graduate Researcher}, American University of Sharjah \hfill \textit{2015 – 2017} \\
\textit{Advisor: Prof.~Ghaleb Husseini}
\begin{tabitemize}[leftmargin=*]
  \item Synthesized estrone-functionalized drug nanocarriers for breast cancer treatment
  \item Optimized ultrasonic parameters for enhanced release control for chemotherapy applications.
  \item Validated drug stability and kinetics with NMR and dynamic light scattering.
  \item Standardized lab protocols; boosted reproducibility and cross-lab collaboration.
  \item Presented at 3 conferences; awarded Best Talk at AUS Biomedical Symposium.
\end{tabitemize}

\section*{Leadership and Community Engagement}
\begin{tabitemize}[leftmargin=*]
  \item \textbf{President}, Early Career Researchers in Polymer Physics (2022–Present): Led 550+ member global network, organized 150+ attendee virtual symposium.
  \item \textbf{Founder \& President}, USF Postdoctoral Scholar Association (2023–Present): Launched NPA-funded ELEVATE Talk Series and DEI programs for 200+ postdocs.
  \item \textbf{Founder \& President}, BYU Chem.~Eng.~Graduate Council (2019–2022): Shaped department policies and spearheaded outreach and recruitment.
  \item \textbf{Chair}, Gordon Research Seminar on Polymer Physics (2026): coordinating speaker selection, session themes, and early-career networking events.
  \item Designed and launched the NPA-funded ELEVATE talk series (5+ events/year) on publishing, inclusive leadership, AI, and career development for 200+ postdocs and graduate researchers.
\end{tabitemize}

%%%%%%%%%%%%%%%%%%%%%%%%%%%%%%%%
%% Articles and Presentations %%
%%%%%%%%%%%%%%%%%%%%%%%%%%%%%%%%
%recompile from start if numbers aren't right
\begin{refsection}[articles]
  \nocite{*}
  \setlength\bibitemsep{0pt}
  \printbibliography[resetnumbers=true,type=article,title={Selected Peer-Reviewed Publications},heading=fix]
\end{refsection}
\noindent\textit{Full list of publications and presentations available at \href{https://linktr.ee/pkawak}{linktr.ee/pkawak}}

\section*{Education}
\textbf{Ph.D.} in Chemical Engineering, Brigham Young University \hfill \textit{2022} \\
\textbf{M.S.} in Chemical Engineering, American University of Sharjah \hfill \textit{2017} \\
\textbf{B.S.} in Chemical Engineering (Econ.~Minor), American University of Sharjah \hfill \textit{2015}

%\vspace{-0.3\baselineskip}

\end{document}
